\section{Programs and checks}
\label{app:code}

The proofs in these notes do not depend on computation. The programs make
the examples checkable and catch mistakes. They are in the directory
\texttt{code/} of the repository
\href{https://github.com/boonsuan/hopcroftkarp}{github.com/boonsuan/hopcroftkarp},
and need only Python~3.10 or later and its standard library (NumPy is used
only by the longest check).

\paragraph{Simulators.} \texttt{algorithm\_h.py} is a literal transcription
of Algorithm~H: it builds the edge nodes as in answer~6, runs step~\Hs1 as in
answer~7, and runs steps \Hs2--\Hs{17} under an explicit program counter. It
records every queue, \A{MARKED} order, \A{BLINK} list, SAP, and the counts
$W_5$, $W_{13}$, $W_{11}$. \texttt{independent\_audit.py} implements the same
rules with ordinary loops and lists and shares no code with it; the two must
produce identical records. \texttt{augmenting.py} contains computations that
do not use Algorithm~H at all: maximum matchings by plain augmenting paths, an
exhaustive list of augmenting paths for small graphs, and $\nu$ by a max-flow
computation.

\paragraph{Examples and figures.} \texttt{examples.py} builds every named
graph of these notes (the staircase graphs $\Gam_n$ and their padding to any
$s$, Knuth's 16-pair graph with the input of Appendix~\ref{app:knuth16}, the
8-pair tree, the halving, caterpillar, six-pair, and triangular chain graphs,
and the catalog of Appendix~\ref{app:catalog}), together with the schedule
and the numbers $T(s)$, $R(s)$, $Q_n$, $\Ch_n$. \texttt{make\_figures.py}
draws every figure that shows a matching, dag, or SAP of a specific graph
by running the simulator, so the pictures show exactly what Algorithm~H does.

\paragraph{Search.} \texttt{search.py} is the search of
Section~\ref{sec:catalog}. For each even $s$ it builds an ordered graph one
round of the Tchoukaillon schedule at a time, rerunning Algorithm~H
from~\Hs1 after every proposed change, and replays what it finds with both
simulators. Its random choices are seeded, so every run finds the same
graphs; they are written to \texttt{code/data/search\_results.json}. The
module's documentation describes the method in full.

\paragraph{Checks.} \texttt{verify.py} runs fifteen named checks, one for
each group of computed statements in these notes, and prints a pass/fail line
for each. They include the hook description of Lemma~\ref{lem:hooks} and the
search of Lemmas~\ref{lem:bfs} and~\ref{lem:dfs} for $n\le60$, the count
$2^m-1$ of SAPs (by exhaustive enumeration for $n\le9$), the counts of
Theorem~W, the padding to every size, the claim about the corner edge in
Section~\ref{sec:gamma}, the values of $\Ch_n$ and the comparison with $Q_n$,
the replay of every example and table, $\nu=f$ in the catalog and in the
graphs found by \texttt{search.py}, the tree of
Section~\ref{sec:no-cover}, and the six-pair and chain obstructions.

\paragraph{Reproducing.} From \texttt{code/}:
\begin{center}
\begin{tabular}{ll}
\texttt{python verify.py} & all checks, about 11 seconds\\
\texttt{python verify.py -{}-full} & full ranges, including $\Ch_n$ for $n\le10^5$\\
& and all 103680 six-pair inputs (about 40 seconds)\\
\texttt{python make\_figures.py} & redraws \texttt{report/figures/gen\_*}\\
\texttt{python search.py 46} & searches every even $s\le46$ (about two minutes)\\
\texttt{python algorithm\_h.py gamma 5} & prints a run (try \texttt{-{}-help})
\end{tabular}
\end{center}
and from \texttt{report/}, two runs of \texttt{pdflatex main.tex} build these
notes. The website \href{https://hopcroftkarp.site}{hopcroftkarp.site} has its own port of the
simulator, tested step by step against traces exported from
\texttt{algorithm\_h.py}.
